\documentclass[a4paper, 11pt]{jsarticle}
\usepackage{amsmath, amssymb}
\usepackage[top=20mm, bottom=30mm, left=20mm, right=20mm, footskip=10mm]{geometry}
\usepackage{multicol}
\usepackage{enumitem}

% ==========================================
% 2段組みのレイアウト調整
% ==========================================
\setlength{\columnseprule}{0.4pt}
\setlength{\columnsep}{1.5cm}

\begin{document}
\pagestyle{empty}

% --- ヘッダー部分 ---
\begin{center}
    {\Large \textbf{二次関数\quad グラフと$x$軸の共有点を求める問題}}
\end{center}
\begin{flushright}
    年\hspace{1cm}組\hspace{1cm}番\hspace{1cm}氏名\hspace*{3.5cm} % ←ここを3.5cmに変更しました
\end{flushright}
\vspace{3mm}

% --- 問題の指示文 ---
\noindent\textbf{問題：} 次の二次関数のグラフと $x$ 軸の共有点の座標を求めよ。
\vspace{5mm}

% ==========================================
% multicols* でページ下端（bottom=30mmのライン）まで伸ばす
% ==========================================
\begin{multicols*}{2}
\begin{enumerate}[label=(\arabic*)]
    % --- 左列（1〜5問目） ---
    \item $y = x^2 - 4x + 3$
    \vspace*{\fill}

    \item $y = x^2 + 2x - 8$
    \vspace*{\fill}

    \item $y = 2x^2 - 5x + 2$
    \vspace*{\fill}

    \item $y = x^2 - 6x + 9$
    \vspace*{\fill}

    \item $y = -x^2 + 3x + 4$
    \vspace*{\fill}
    
    \columnbreak

    % --- 右列（6〜10問目） ---
    \item $y = x^2 - 2x - 3$
    \vspace*{\fill}

    \item $y = x^2 - 9$
    \vspace*{\fill}

    \item $y = 3x^2 + 5x - 2$
    \vspace*{\fill}

    \item $y = -x^2 + 4x - 3$
    \vspace*{\fill}

    \item $y = x^2 + 4x + 4$
    \vspace*{\fill}

\end{enumerate}
\end{multicols*}
% ==========================================

\end{document}